\section{Introduction}

In modern AI, individuality is a misconception. Large labs, such as OpenAI and Anthropic, frequently publicise how their models have completed unsolved problems in mathematics, the natural sciences, and coding \parencite{georgiev2025, bubeck2025, elkishky2025}. At first glance, this sounds like a single agent reaching a eureka moment. However, the reality looks much closer to agents coordinating in large groups, sometimes 10,000 or more strong, to reach an insight \parencite{openai2026navierstokes}. This coordination does not detract from the frontier models’ incredible achievements. However, it means that frontier intelligence now pertains to a \textit{collective} rather than an individual.

The latest safety incidents associated with frontier models also have a collective property. A clear example is the recent OpenAI-Hugging Face incident (OAI-HF), where roughly 1,200 AI agents performed a coordinated cybersecurity attack against the two companies \parencite{metr-2026-openai-hugging-face-incident-investigation, openai_hugging_2026}. This agent ‘swarm’ exhibited many of the properties we might associate with a cult or similar fanatically devoted group \parencite{amodei_we_2026}. They followed instructions from a central coordinator who asked devotees to commit crimes based on paranoid and hyperbolic assumptions about what it would take to succeed on a narrowly defined task \parencite{metr-2026-openai-hugging-face-incident-investigation}. The devotees sacrificed themselves to advance the group and, most concerningly, failed to alert a human even when they knew the group's actions were wrong \parencite{metr-2026-openai-hugging-face-incident-investigation}.

Similar, albeit less malign, incidents are occurring elsewhere on the internet, with agents infiltrating obscure websites and using them as message boards to coordinate on tasks \parencite{kleinman_openai_2026}. Indeed, the volume of these events is likely much higher than what has been reported, since we can only spot collectives that have failed to cover their tracks \parencite{motwani2024}. When we combine these empirical observations with projections about the growth of agentic internet traffic \parencite{cloudflare_traffic_nodate}, scenarios with extremely costly socioeconomic ramifications, such as ‘persistent botnets’ \parencite{amodei_we_2026}, now seem plausible. Leaders from frontier labs have recognised this threat and are rallying behind calls to slow progress and ‘pace the frontier’ \parencite{bradshaw_rivals_2026}. Nevertheless, AI safety researchers broadly agree that the field lacks effective techniques for aligning collectives of AI agents \parencite{hammond2025}.

Artificial Societies is uniquely positioned to make theoretical and empirical contributions to aligning artificial collective intelligence (ACI), a term we refer to as group alignment. Alongside authoring the first paper on AI swarms \parencite{he_artificial_2024}, our simulation technology uses large collectives of agentic ‘personas’ to model how groups of people behave in different scenarios. As an AI company with a mission to develop socially intelligent technology, we also have a responsibility to advance safety and minimise the likelihood of adverse socioeconomic effects from AI. This blog post is written in that spirit.

Below, we propose a definition of group alignment that distinguishes the alignment of one agent from the alignment of a collective and explains how agents which are aligned with human values at an individual level can coordinate as a group to subvert us. Following these theoretical contributions, we outline several numerical experiments showing why a method of group alignment we call ‘mixed swarms’ could reduce the risks of agentic collectives. We conclude with some next steps for safety research on group alignment.
